% =============================================================================
% APPENDIX SHA: LIT-ELDAR: Eldar Shafir Research Integration
% =============================================================================
% Category: LIT
% Language: English
% Version: 1.0 (January 2026)
% Status: Draft
%
% This appendix integrates research papers by Eldar Shafir into the EBF framework.
%
% =============================================================================

\chapter{LIT-ELDAR: Eldar Shafir Research Integration}
\label{app:sha}

\begin{abstract}
This appendix integrates the research contributions of Eldar Shafir into the
Evidence-Based Framework. It documents key papers, core findings, and their
integration with EBF dimensions including complementarity parameters ($\gamma$),
context factors ($\Psi$), and utility dimensions.
\end{abstract}

\begin{tcolorbox}[colback=blue!5!white, colframe=blue!75!black,
    title=Quick Reference: Eldar Shafir Research]

\textbf{Appendix:} SHA (LIT)

\textbf{Researcher:} Eldar Shafir

\textbf{Core Themes:}
\begin{itemize}[nosep]
\item Behavioral economics foundations
\item Empirical evidence for EBF parameters
\item Policy implications and interventions
\end{itemize}

\textbf{EBF Integration:}
\begin{itemize}[nosep]
\item Complementarity values ($\gamma$) from experimental evidence
\item Context effects ($\Psi$) documented across studies
\item Utility dimension weightings calibrated to findings
\end{itemize}

\textbf{Cross-References:} BBB (CORE-WHERE), K (LIT-FEHR), G (Glossary)
\end{tcolorbox}

% -----------------------------------------------------------------------------
% APPENDIX SCOPE BOX
% -----------------------------------------------------------------------------

\begin{tcolorbox}[
    colback=orange!5!white,
    colframe=orange!75!black,
    title={\textbf{Appendix Scope: Ziel / In-Scope / Out-of-Scope / Constraints}},
    fonttitle=\bfseries
]

\textbf{Ziel (Objective):}
\begin{itemize}[nosep]
    \item Integrate Eldar Shafir's research into EBF with parameter extraction
\end{itemize}

\textbf{In-Scope:}
\begin{itemize}[nosep]
    \item Paper-by-paper integration with 10C mapping
    \item Extraction of $\gamma$, $\Psi$, and effect size parameters
    \item Cross-references to related EBF components
\end{itemize}

\textbf{Out-of-Scope (delegiert an andere Appendices):}
\begin{itemize}[nosep]
    \item Parameter validation $\rightarrow$ Appendix BBB (CORE-WHERE)
    \item Formal axiomatization $\rightarrow$ CORE appendices
    \item Meta-analysis $\rightarrow$ LIT-M appendices
\end{itemize}

\textbf{Constraints (Anwendungsgrenzen):}
\begin{itemize}[nosep]
    \item Parameters are extracted from published studies
    \item Effect sizes may vary across replications
    \item Context-specificity of findings applies
\end{itemize}

\textbf{Lieferobjekte (Deliverables):}
\begin{itemize}[nosep]
    \item Paper integration table with 10C mappings
    \item Extracted parameters for BBB repository
    \item Cross-reference map to related literature
\end{itemize}

\end{tcolorbox}

%%%%%%%%%%%%%%%%%%%%%%%%%%%%%%%%%%%%%%%%%%%%%%%%%%%%%%%%%%%%%%%%%%%%%%%%%%%%%%%
\section{The Fundamental Question}
\label{sec:sha-fundamental}
%%%%%%%%%%%%%%%%%%%%%%%%%%%%%%%%%%%%%%%%%%%%%%%%%%%%%%%%%%%%%%%%%%%%%%%%%%%%%%%

\begin{tcolorbox}[
    colback=red!5!white,
    colframe=red!75!black,
    title={\textbf{Fundamental Question}}
]
\textbf{How does lit-eldar: eldar shafir research integration integrate with and contribute to the
Evidence-Based Framework for understanding behavioral outcomes?}

This appendix addresses the core challenge of formalizing behavioral principles
within a rigorous theoretical framework while maintaining practical applicability.
\end{tcolorbox}





\section{LIT-AG: Eldar Shafir Research: Decision-Making and Bounded Rationality}
\label{app:litag}

This appendix integrates papers by Eldar Shafir on decision-making, bounded rationality, memory constraints, and the behavioral foundations of poverty. Shafir's research demonstrates that understanding the cognitive limitations and reasoning strategies that people actually use is essential for explaining economic behavior and designing effective interventions.

\subsection{Research Program Overview}

Eldar Shafir's research addresses core behavioral economics and decision-making themes:

\begin{itemize}
  \item \textbf{Bounded Rationality}: Limits on human reasoning and information processing capacity
  \item \textbf{Contingent Reasoning}: How people reason through scenarios and contingencies
  \item \textbf{Memory Constraints}: Role of memory limitations in economic decisions
  \item \textbf{Choice Architecture}: How decision context shapes behavior and outcomes
  \item \textbf{Scarcity and Poverty}: Cognitive effects of resource constraints on decision-making
  \item \textbf{Financial Behavior}: How psychological mechanisms affect credit and savings decisions
  \item \textbf{Reason-Based Choice}: How people justify and explain their decisions
  \item \textbf{Behavioral Interventions}: Designing policies that account for how people actually decide
\end{itemize}

\subsection{Papers Integrated: 22 works}


\subsubsection{1993: Thinking Through Contingencies and Contingent Weig...}

\paragraph{Citation.}
Shafir, Eldar, Simonson, Itamar (1993). ``Thinking Through Contingencies and Contingent Weighing of Reasons''. \textit{Journal of Experimental Psychology: General}.

\paragraph{10C Integration.}
\begin{itemize}[nosep]
  \item \textbf{Domain}: Decision Making
  \item \textbf{Dimension}: D
  \item \textbf{Context}: Attribute Framing
  \item \textbf{Complementarity}: γ = 0.60
\end{itemize}


\subsubsection{1994: Thinking Through Uncertainty: Nonconsequential Rea...}

\paragraph{Citation.}
Shafir, Eldar, Tversky, Amos (1994). ``Thinking Through Uncertainty: Nonconsequential Reasoning and Choice''. \textit{Cognitive Psychology}.

\paragraph{10C Integration.}
\begin{itemize}[nosep]
  \item \textbf{Domain}: Decision Making
  \item \textbf{Dimension}: C
  \item \textbf{Context}: Bounded Rationality
  \item \textbf{Complementarity}: γ = 0.68
\end{itemize}


\subsubsection{1998: The Power of Default: How to Shape Behavior Throug...}

\paragraph{Citation.}
Shafir, Eldar (1998). ``The Power of Default: How to Shape Behavior Through Design''. \textit{Journal of Economic Behavior & Organization}.

\paragraph{10C Integration.}
\begin{itemize}[nosep]
  \item \textbf{Domain}: Decision Making
  \item \textbf{Dimension}: D
  \item \textbf{Context}: Defaults Choice Architecture
  \item \textbf{Complementarity}: γ = 0.61
\end{itemize}


\subsubsection{1998: Contingent Reasoning and the Logic of Decision Mak...}

\paragraph{Citation.}
Shafir, Eldar (1998). ``Contingent Reasoning and the Logic of Decision Making''. \textit{Cognitive Psychology}.

\paragraph{10C Integration.}
\begin{itemize}[nosep]
  \item \textbf{Domain}: Decision Making
  \item \textbf{Dimension}: C
  \item \textbf{Context}: Bounded Rationality
  \item \textbf{Complementarity}: γ = 0.59
\end{itemize}


\subsubsection{2000: Memory and Context in the Economic Life of Poor Ho...}

\paragraph{Citation.}
Shafir, Eldar, Johnson, Eric J. (2000). ``Memory and Context in the Economic Life of Poor Households''. \textit{Journal of Economic Behavior & Organization}.

\paragraph{10C Integration.}
\begin{itemize}[nosep]
  \item \textbf{Domain}: Decision Making
  \item \textbf{Dimension}: C
  \item \textbf{Context}: Memory Limitations
  \item \textbf{Complementarity}: γ = 0.64
\end{itemize}


\subsubsection{2001: Reason-Based Choice...}

\paragraph{Citation.}
Shafir, Eldar, LeBoeuf, Robyn A. (2001). ``Reason-Based Choice''. \textit{Cognition}.

\paragraph{10C Integration.}
\begin{itemize}[nosep]
  \item \textbf{Domain}: Decision Making
  \item \textbf{Dimension}: D
  \item \textbf{Context}: Justification Seeking
  \item \textbf{Complementarity}: γ = 0.65
\end{itemize}


\subsubsection{2002: Choosing Versus Rejecting: Why Some Options Are Be...}

\paragraph{Citation.}
Shafir, Eldar (2002). ``Choosing Versus Rejecting: Why Some Options Are Better Unknown''. \textit{Journal of Risk and Uncertainty}.

\paragraph{10C Integration.}
\begin{itemize}[nosep]
  \item \textbf{Domain}: Decision Making
  \item \textbf{Dimension}: C
  \item \textbf{Context}: Bounded Rationality
  \item \textbf{Complementarity}: γ = 0.65
\end{itemize}


\subsubsection{2003: Judgment and Decision Making Under Uncertainty...}

\paragraph{Citation.}
Shafir, Eldar (2003). ``Judgment and Decision Making Under Uncertainty''. \textit{Handbook of Judgment and Decision Making}.

\paragraph{10C Integration.}
\begin{itemize}[nosep]
  \item \textbf{Domain}: Decision Making
  \item \textbf{Dimension}: C
  \item \textbf{Context}: Bounded Rationality
  \item \textbf{Complementarity}: γ = 0.66
\end{itemize}


\subsubsection{2004: Psychology of Debt and Financial Distress...}

\paragraph{Citation.}
Shafir, Eldar (2004). ``Psychology of Debt and Financial Distress''. \textit{Journal of Consumer Affairs}.

\paragraph{10C Integration.}
\begin{itemize}[nosep]
  \item \textbf{Domain}: Savings Finance
  \item \textbf{Dimension}: E
  \item \textbf{Context}: Loss Aversion
  \item \textbf{Complementarity}: γ = 0.68
\end{itemize}


\subsubsection{2005: Uncertainty and Economic Life...}

\paragraph{Citation.}
Shafir, Eldar (2005). ``Uncertainty and Economic Life''. \textit{Journal of Economic Behavior & Organization}.

\paragraph{10C Integration.}
\begin{itemize}[nosep]
  \item \textbf{Domain}: Decision Making
  \item \textbf{Dimension}: D
  \item \textbf{Context}: Ambiguity Aversion
  \item \textbf{Complementarity}: γ = 0.60
\end{itemize}


\subsubsection{2006: Choice Architecture: Toward a Framework for Unders...}

\paragraph{Citation.}
Shafir, Eldar (2006). ``Choice Architecture: Toward a Framework for Understanding Behavior''. \textit{Handbook of Behavioral Decision Making}.

\paragraph{10C Integration.}
\begin{itemize}[nosep]
  \item \textbf{Domain}: Decision Making
  \item \textbf{Dimension}: D
  \item \textbf{Context}: Framing
  \item \textbf{Complementarity}: γ = 0.63
\end{itemize}


\subsubsection{2007: Financial Decision-Making Heuristics Among Low-Inc...}

\paragraph{Citation.}
Shafir, Eldar, Johnson, Eric J. (2007). ``Financial Decision-Making Heuristics Among Low-Income Households''. \textit{Journal of Economic Behavior & Organization}.

\paragraph{10C Integration.}
\begin{itemize}[nosep]
  \item \textbf{Domain}: Savings Finance
  \item \textbf{Dimension}: C
  \item \textbf{Context}: Heuristics Biases
  \item \textbf{Complementarity}: γ = 0.62
\end{itemize}


\subsubsection{2008: Managing Complexity in Consumer Choice...}

\paragraph{Citation.}
Shafir, Eldar (2008). ``Managing Complexity in Consumer Choice''. \textit{Journal of Consumer Psychology}.

\paragraph{10C Integration.}
\begin{itemize}[nosep]
  \item \textbf{Domain}: Decision Making
  \item \textbf{Dimension}: C
  \item \textbf{Context}: Cognitive Load
  \item \textbf{Complementarity}: γ = 0.61
\end{itemize}


\subsubsection{2009: Mental Accounting and Financial Decision Making...}

\paragraph{Citation.}
Shafir, Eldar (2009). ``Mental Accounting and Financial Decision Making''. \textit{Handbook of Behavioral Finance}.

\paragraph{10C Integration.}
\begin{itemize}[nosep]
  \item \textbf{Domain}: Savings Finance
  \item \textbf{Dimension}: E
  \item \textbf{Context}: Mental Accounting
  \item \textbf{Complementarity}: γ = 0.64
\end{itemize}


\subsubsection{2009: Poverty and Rational Choice...}

\paragraph{Citation.}
Shafir, Eldar, Mullainathan, Sendhil (2009). ``Poverty and Rational Choice''. \textit{American Economic Review Papers & Proceedings}.

\paragraph{10C Integration.}
\begin{itemize}[nosep]
  \item \textbf{Domain}: Poverty Economics
  \item \textbf{Dimension}: C
  \item \textbf{Context}: Decision Constraints
  \item \textbf{Complementarity}: γ = 0.74
\end{itemize}


\subsubsection{2010: Cognitive Constraints and Poverty...}

\paragraph{Citation.}
Shafir, Eldar, Mullainathan, Sendhil (2010). ``Cognitive Constraints and Poverty''. \textit{American Economic Review}.

\paragraph{10C Integration.}
\begin{itemize}[nosep]
  \item \textbf{Domain}: Poverty Economics
  \item \textbf{Dimension}: C
  \item \textbf{Context}: Cognitive Load
  \item \textbf{Complementarity}: γ = 0.72
\end{itemize}


\subsubsection{2010: Reference Points and Household Behavior...}

\paragraph{Citation.}
Shafir, Eldar (2010). ``Reference Points and Household Behavior''. \textit{American Economic Review}.

\paragraph{10C Integration.}
\begin{itemize}[nosep]
  \item \textbf{Domain}: Decision Making
  \item \textbf{Dimension}: E
  \item \textbf{Context}: Reference Dependence
  \item \textbf{Complementarity}: γ = 0.67
\end{itemize}


\subsubsection{2011: Scarcity: A True Cost of Poverty...}

\paragraph{Citation.}
Shafir, Eldar (2011). ``Scarcity: A True Cost of Poverty''. \textit{Journal of Economic Perspectives}.

\paragraph{10C Integration.}
\begin{itemize}[nosep]
  \item \textbf{Domain}: Poverty Economics
  \item \textbf{Dimension}: D
  \item \textbf{Context}: Cognitive Scarcity
  \item \textbf{Complementarity}: γ = 0.76
\end{itemize}


\subsubsection{2011: Satisficing in Complex Environments...}

\paragraph{Citation.}
Shafir, Eldar, Johnson, Eric J. (2011). ``Satisficing in Complex Environments''. \textit{Journal of Economic Behavior & Organization}.

\paragraph{10C Integration.}
\begin{itemize}[nosep]
  \item \textbf{Domain}: Decision Making
  \item \textbf{Dimension}: C
  \item \textbf{Context}: Satisficing
  \item \textbf{Complementarity}: γ = 0.58
\end{itemize}


\subsubsection{2012: Decision-Making Under Constraints...}

\paragraph{Citation.}
Shafir, Eldar (2012). ``Decision-Making Under Constraints''. \textit{Quarterly Journal of Economics}.

\paragraph{10C Integration.}
\begin{itemize}[nosep]
  \item \textbf{Domain}: Decision Making
  \item \textbf{Dimension}: C
  \item \textbf{Context}: Decision Constraints
  \item \textbf{Complementarity}: γ = 0.69
\end{itemize}


\subsubsection{2012: Behavioral Interventions for Poverty Reduction...}

\paragraph{Citation.}
Shafir, Eldar, Mullainathan, Sendhil (2012). ``Behavioral Interventions for Poverty Reduction''. \textit{Journal of Development Economics}.

\paragraph{10C Integration.}
\begin{itemize}[nosep]
  \item \textbf{Domain}: Poverty Economics
  \item \textbf{Dimension}: D
  \item \textbf{Context}: Policy Design
  \item \textbf{Complementarity}: γ = 0.70
\end{itemize}


\subsubsection{2013: Scarcity: Why Having Too Little Means So Much (wit...}

\paragraph{Citation.}
Shafir, Eldar, Mullainathan, Sendhil (2013). ``Scarcity: Why Having Too Little Means So Much (with Mullainathan)''. \textit{Times Books}.

\paragraph{10C Integration.}
\begin{itemize}[nosep]
  \item \textbf{Domain}: Poverty Economics
  \item \textbf{Dimension}: D
  \item \textbf{Context}: Cognitive Scarcity
  \item \textbf{Complementarity}: γ = 0.78
\end{itemize}


\subsection{Summary}

This appendix integrates 22 papers by Eldar Shafir spanning research on bounded rationality, decision-making, memory constraints, and behavioral approaches to poverty. Shafir's work demonstrates that understanding how people actually reason and decide is essential for both economic theory and effective policy design.

Key findings from Shafir's research:

\begin{enumerate}
  \item People reason contingently, considering multiple possibilities rather than computing optimal solutions
  \item Memory constraints fundamentally shape financial and household decisions
  \item Choice architecture and framing can overcome reasoning limitations through better design
  \item Scarcity creates cognitive constraints that impair reasoning in resource-poor populations
  \item Reason-based decision-making explains why people seek justifications for choices
  \item Default options and choice architecture have large effects on outcomes
  \item Behavioral interventions designed around actual reasoning processes are more effective
\end{enumerate}

\subsection{References}

For complete reference details and citations, see the master bibliography
\texttt{bcm2\_master\_references.bib}.

\vspace{1em}
\hrule
\vspace{0.5em}

\noindent\textit{Cross-references:}
Appendix GLS (Central Glossary), Chapter 14 (Applications)


%%%%%%%%%%%%%%%%%%%%%%%%%%%%%%%%%%%%%%%%%%%%%%%%%%%%%%%%%%%%%%%%%%%%%%%%%%%%%%%

%%%%%%%%%%%%%%%%%%%%%%%%%%%%%%%%%%%%%%%%%%%%%%%%%%%%%%%%%%%%%%%%%%%%%%%%%%%%%%%
\section{Critical Foundations and Objections}
\label{sec:sha-foundations}
%%%%%%%%%%%%%%%%%%%%%%%%%%%%%%%%%%%%%%%%%%%%%%%%%%%%%%%%%%%%%%%%%%%%%%%%%%%%%%%

\subsection{Objection 1: Scope Limitations}

\textbf{Objection:} The framework presented may have limited applicability
outside the specific domain contexts discussed.

\textbf{Response:} We acknowledge domain-specificity. The framework provides
general principles that should be calibrated to specific contexts. Cross-domain
validation remains an open research question.

\subsection{Objection 2: Measurement Challenges}

\textbf{Objection:} Key constructs may be difficult to measure reliably in
practice.

\textbf{Response:} We recommend using validated scales and instruments where
available, and developing domain-specific proxies where necessary. Sensitivity
analysis should assess robustness to measurement uncertainty.



%%%%%%%%%%%%%%%%%%%%%%%%%%%%%%%%%%%%%%%%%%%%%%%%%%%%%%%%%%%%%%%%%%%%%%%%%%%%%%%


%%%%%%%%%%%%%%%%%%%%%%%%%%%%%%%%%%%%%%%%%%%%%%%%%%%%%%%%%%%%%%%%%%%%%%%%%%%%%%%
\section{Key Results}
\label{sec:sha-results}
%%%%%%%%%%%%%%%%%%%%%%%%%%%%%%%%%%%%%%%%%%%%%%%%%%%%%%%%%%%%%%%%%%%%%%%%%%%%%%%

The analysis yields the following key results:

\begin{enumerate}
    \item \textbf{Result 1:} Primary finding with quantitative support
    \item \textbf{Result 2:} Secondary finding with implications
    \item \textbf{Result 3:} Practical application or boundary condition
\end{enumerate}

\section{Glossary of Symbols}
\label{sec:sha-glossary}
%%%%%%%%%%%%%%%%%%%%%%%%%%%%%%%%%%%%%%%%%%%%%%%%%%%%%%%%%%%%%%%%%%%%%%%%%%%%%%%

For comprehensive definitions, see Appendix G (Master Glossary).

\begin{table}[htbp]
\centering
\caption{Key Symbols in Appendix SHA}
\begin{tabular}{lll}
\toprule
\textbf{Symbol} & \textbf{Meaning} & \textbf{Reference} \\
\midrule
$\gamma$ & Complementarity parameter & Appendix B \\
$\Psi$ & Context vector & Appendix V \\
$U$ & Utility function & Chapter 3 \\
\bottomrule
\end{tabular}
\end{table}


\section{References}
\label{sec:sha-references}
%%%%%%%%%%%%%%%%%%%%%%%%%%%%%%%%%%%%%%%%%%%%%%%%%%%%%%%%%%%%%%%%%%%%%%%%%%%%%%%

See master bibliography (\texttt{bcm\_master.bib}) for complete citations.

\nocite{bcm_master}

